\documentclass[11pt,a4paper]{article}

\usepackage[margin=2.2cm]{geometry}

\usepackage[T1]{fontenc}\usepackage{lmodern,booktabs,graphicx,microtype}

\setlength{\parindent}{0pt}\setlength{\parskip}{0.65em}

\begin{document}

{\LARGE\bfseries Iranian REIT Management Brief}\par

Prepared 08 October 2026\par

\textbf{Purpose.} Assess whether the observed REIT returns delivered useful housing exposure and preserved purchasing power over the two-year review period.

\textbf{Main finding.} Kelid and Danik substantially underperformed the housing proxy. Arzesh Maskan finished close to it. All three underperformed USD/IRR and TEDPIX. Monthly share returns showed little co-movement with the housing proxy, and the two portfolios with allocation data held substantial non-property assets.

\textbf{Investor performance.} Arzesh Maskan returned 156.8\%, Kelid 97.9\%, and Danik 74.3\%, compared with 160.7\% for housing, 267.2\% for TEDPIX, and 312.8\% for USD/IRR. The indicative housing shortfalls were 3.9, 62.8, and 86.4 percentage points, respectively. Measured against USD/IRR, ending purchasing power fell by 37.8\%, 52.1\%, and 57.8\% respectively. Table 1 distinguishes nominal gains from benchmark shortfalls.

\textbf{Housing exposure.} Arzesh's direct-property allocation declined from 83.3\% to 68.7\%; its latest fixed-income share was 26.7\%. Danik's direct-property allocation declined from 89.6\% to 60.2\%, but averaged only 43.1\% across the 24 snapshots and reached a low of 23.9\%. Its fixed-income share reached 75.4\% during the period. At the latest month end, Danik held 12.9\% in fixed income and 26.9\% in other assets, primarily other/receivables. These portfolios did not provide pure direct-property exposure. Table 2 summarizes the observed composition.

\textbf{Housing tracking.} Over 24 months, housing-return correlations were $-0.064$ for Arzesh, $+0.094$ for Kelid, and $-0.178$ for Danik, based on 21--22 paired monthly returns. These coefficients indicate weak contemporaneous co-movement with this benchmark. They do not measure cumulative underperformance or prove that the underlying properties failed to appreciate. Table 3 reports the sample sizes.

\textbf{Management implication.} The evidence warrants a review of portfolio composition and shareholder-return delivery. It does not establish mismanagement as the cause. Attribution requires complete distributions, historical NAV and share-price discounts, property operating results, fees, financing costs, and a holdings-based return bridge. No allocation conclusion can be drawn for Kelid from the supplied workbook, which covers Arzesh and Danik only.

\textbf{Essential qualifications.} Market comparisons run from 4 October 2024 to 2 October 2026. Housing runs from 21 October 2024 to 22 September 2026 and is a monthly Kilid listing-price proxy excluding rental income and ownership costs; its return gap is therefore indicative. REIT returns assume immediate reinvestment of only three recorded approved distributions; payout completeness and actual cash dates remain unresolved. The allocation workbook is user-supplied and contains three flagged reconstructed observations. Kakh is excluded from this two-year performance assessment because its history is too short.

\newpage

{\Large\bfseries Supporting tables}\par

\textbf{Table 1 Investor performance and benchmark gaps}

\begin{center}
\small
\resizebox{\linewidth}{!}{%
\begin{tabular}{lrrr}
\toprule
Asset & Cumulative return & Housing gap pp & USD-relative wealth \\
\midrule
Arzesh Maskan & +156.8\% & -3.9 & -37.8\% \\
Kelid & +97.9\% & -62.8 & -52.1\% \\
Danik & +74.3\% & -86.4 & -57.8\% \\
Tehran housing & +160.7\% & -- & -- \\
TEDPIX & +267.2\% & -- & -- \\
USD/IRR & +312.8\% & -- & -- \\
\bottomrule
\end{tabular}%
}
\end{center}

{\small pp means percentage points. USD-relative wealth equals $(1+R_{fund})/(1+R_{USD})-1$. It measures the change in dollar purchasing power, not the arithmetic return difference. Housing uses the different monthly dates stated above.}\par

\textbf{Table 2 Arzesh and Danik portfolio composition}

\begin{center}
\small
\resizebox{\linewidth}{!}{%
\begin{tabular}{lrrrrr}
\toprule
Fund & First housing & Latest housing & Mean housing & Latest fixed income & Latest other \\
\midrule
Arzesh Maskan & 83.3\% & 68.7\% & 76.5\% & 26.7\% & 4.6\% \\
Danik & 89.6\% & 60.2\% & 43.1\% & 12.9\% & 26.9\% \\
\bottomrule
\end{tabular}%
}
\end{center}

{\small Coverage is 1403/07--1405/06, with 24 snapshots per fund. Mean housing is the equally weighted mean of month-end accounting shares. Other includes cash, equity, and other/receivables. Latest housing, fixed income, and other sum to 100\% before rounding.}\par

\textbf{Table 3 Monthly housing return correlation over 24 months}

\begin{center}
\small
\resizebox{\linewidth}{!}{%
\begin{tabular}{lrr}
\toprule
Fund & Pearson correlation & Paired months \\
\midrule
Arzesh Maskan & -0.064 & 22 \\
Kelid & +0.094 & 21 \\
Danik & -0.178 & 22 \\
\bottomrule
\end{tabular}%
}
\end{center}

{\small Same-month comparisons use reconstructed REIT returns and the housing proxy. Missing returns are not filled. All three rows meet the minimum coverage rule.}\par

\textbf{Source.} Existing processed research inputs underlying \texttt{REIT\_CROSS\_ASSET\_ANALYSIS.tex}, including the user-supplied Codal-derived asset-mix workbook. This brief is a summary of that report; no market-source refresh was performed.

\end{document}
